\section{A Timeout Exchanger}

In previous chapters, we saw a couple of implementations of an exchanger: an
object that allows two threads to swap values of some type~|A|.  In this
section we consider an exchanger with a timeout: if a thread fails to exchange
within some period of time, it will timeout and fail.  We will also allow the
exchange to fail under another couple of circumstances, which we discuss below.

In order to capture the possibility of failure, we give the |exchange|
operation the following signature:
%
\begin{scala}
  def exchange(x: A): Option[A]
\end{scala}
% 
A return value of |Some(y)| will indicate that the thread has successfully
exchanged with another thread and obtained~|y|; a return value of |None| will
indicate that the exchange failed. 

In order to implement the timeout, we will need another form of the |wait|
operation.  The operation |wait(delay)| is like a standard |wait()|, except if
it doesn't receive a signal, it will timeout after |delay| milliseconds (ms).
In each case, the thread needs to then wait to reobtain the lock on the
monitor.  The client code is responsible for working out whether it received a
signal or not.

There is another version, |wait(millis, nanos)|, which waits for |millis|
milliseconds and |nanos| nanoseconds (ns); equivalently, it waits for
$1,000,000 \times \sm{millis} + \sm{nanos}$\,ns.  

A millisecond is rather a long time for a thread to wait for another thread.
(Delays measured in milliseconds might be appropriate to wait for a human or
for a network communication.)  We will therefore use this latter form, and
measure the timeout time for the exchanger in nanoseconds.  We will use an
|Int| variable |delay| to record the delay, so the maximum delay will be
|Int.MaxValue| $= 2^{31}-1$\, ns, about two seconds.

Our implementation is in Figure~\ref{fig:TimeoutExchanger}.  It is somewhat
similar to the exchanger in Figure~\ref{fig:MonitorExchanger}.  It uses a
variable |stage| that records which stage has been reached.  This variable
will store one of the following values.
%
\begin{description}
\item[{\scalastyle Empty}:] No thread is currently exchanging;

\item[{\scalashape Filled}:] The first thread has deposited its value;

\item[{\scalashape Exchanging}:] An exchange is underway: the second thread
  has taken the first thread's value, and swapped it with its own value.
\end{description}


%%%%%

\begin{figure}
\begin{scala}
class TimeoutExchanger[A](delay: Int){
  private var slot: A = _

  private val Empty = 0; private val Filled = 1; private val Exchanging = 2

  private var stage: Int = Empty

  def exchange(x: A): Option[A] = synchronized{
    if(stage == Empty){
      // Deposit my value and wait.
      slot = x; stage = Filled; wait(0, delay)
      if(stage == Exchanging){ stage = Empty; Some(slot) } // Success.
      else{ assert(stage == Filled); stage = Empty; None } // No joy.
    }
    else if(stage == Filled){  // Take value in £slot£; deposit my value there.
      val result = slot; slot = x; stage = Exchanging; notify(); Some(result)
    }
    else{ assert(stage == Exchanging); None } // Another exchange is under way.
  }
}
\end{scala}
\caption{A timeout exchanger.}
\label{fig:TimeoutExchanger}
\end{figure}

%%%%%

If a thread finds that $\sm{stage} = \sm{Empty}$, it is the first thread in
the exchange.  It stores its value, sets $\sm{stage} = \sm{Filled}$, and
waits for up to |delay|\,ns.  When it wakes up, if $\sm{stage} =
\sm{Exchanging}$, another thread has deposited a value, so it can return that
value, resetting for the next round.  In fact, it acts the same way if it
times out just before another thread deposits its value; this behaviour is
desirable.  Alternatively, if it finds that |stage| is still |Filled| when it
wakes up, it has failed to exchange, so resets for the next stage, and returns
|None|.  This latter behaviour would also result from a spurious wakeup; we
consider this behaviour acceptable: to detect the spurious wakeup would
require comparing the system time before and after the |wait|, which is
awkward.  

If the thread initially finds that $\sm{stage} = \sm{Filled}$, it is the
second thread in the exchange.  It swaps its value with the value stored by
the first thread, sets |stage| to |Exchanging|, signals to the first thread,
and returns the first thread's value. 

If the thread initially finds that $\sm{stage} = \sm{Exchanging}$, it means
that two other threads are currently exchanging.  Clearly, this thread should
not interfere with that exchange.  We make the decision to treat this case as
a failure to exchange, and just return |None|.  An alternative would be to
wait while $\sm{stage} = \sm{Exchanging}$, for up to |delay|\,ns.  If
subsequently this thread becomes the first thread in an exchange, the length
of its subsequent wait would have to be adjusted based on the length of the
earlier wait.  This is rather messy, so we don't adopt this approach. 

\begin{instruction}
Study the details of the implementation.
\end{instruction}

The timeout exchanger can be tested much the same as the basic exchanger in
Section~\ref{sec:exchanger}, except we should also allow an attempted exchange
to fail.  Note that we cannot require an attempted exchange to succeed, even
if it overlaps in time with another attempt: it is possible that unfortunate
scheduling means that one of other is delayed, so as to prevent the exchange.
